\chapter{Aortic Valve Replacement}
\index{Aortic Valve Replacement}

\section*{Definition and Overview}
Aortic valve replacement is an operation to remove a diseased aortic valve and replace it with a prosthetic (artificial) valve, restoring normal blood flow from the heart to the body. It is the definitive treatment for severe aortic stenosis and for certain cases of aortic regurgitation, endocarditis, and congenital valve disease.

The replacement is performed in two main ways. \textbf{Surgical aortic valve replacement} (SAVR) is open-heart surgery in which the chest is opened, the heart is placed on a bypass machine, and the diseased valve is excised and replaced. \textbf{Transcatheter aortic valve implantation} (TAVI) is a less invasive approach in which a new tissue valve is folded inside a catheter and expanded in place within the old valve, usually through the groin artery. The choice between the approaches depends on age, the type and severity of valve disease, other health conditions, and the anatomy of the heart and arteries.

\section*{Epidemiology}
Aortic valve replacement is one of the most common cardiac operations in the world, and the number performed continues to rise as populations age and age-related calcific aortic stenosis becomes more prevalent. SAVR has been performed for decades and remains well suited to younger, lower-risk patients, offering excellent durability.

TAVI has grown rapidly since its introduction in the 2000s and is now the dominant approach in older adults and in people considered too frail or high-risk for open surgery. As the technique has matured, its use has extended to intermediate-risk and, increasingly, lower-risk younger patients. Rates of both procedures are expected to keep growing, driven by an ageing population and widening indications.

\section*{Aetiology and Pathophysiology}
The valve is replaced because disease has made it unable to perform its normal function:

\begin{itemize}
    \item \textbf{Severe aortic stenosis:} progressive calcification narrows the valve and obstructs blood flow -- by far the most common reason for replacement
    \item \textbf{Aortic regurgitation:} a valve that fails to close properly, allowing blood to leak back into the left ventricle
    \item \textbf{Infective endocarditis:} severe infection that destroys the valve leaflets
    \item \textbf{Bicuspid aortic valve:} a congenital two-leaflet valve that degenerates and calcifies decades earlier than a normal valve
    \item \textbf{Failed previous replacement:} a prosthetic valve that has worn out, degenerated, or become infected, requiring reoperation
    \item \textbf{Mechanism:} in SAVR the diseased valve is removed and a mechanical or tissue valve is sewn into the annulus; in TAVI a tissue valve is compressed onto a stent and expanded within the old valve, which acts as a scaffold
\end{itemize}

\section*{Clinical Features}
Before surgery, most people have symptoms of severe valve disease, such as exertional breathlessness, chest pain, fainting, or fatigue. The operation itself is followed by a characteristic recovery course:

\begin{itemize}
    \item Soreness around the chest wound (after SAVR) or the groin (after TAVI), which settles over days to weeks
    \item Fatigue and reduced energy, usually lasting several weeks
    \item Reduced appetite and some weight loss early in recovery
    \item Mild swelling around the wound site that gradually resolves
    \item Breathlessness and a cough that improve as fluid balance returns to normal
    \item Common early heart rhythm changes, such as atrial fibrillation, which may settle spontaneously or require treatment
    \item A gradual return of energy and exercise capacity over weeks to months, aided by cardiac rehabilitation
\end{itemize}

\section*{Differential Diagnosis}
The main clinical decisions are not which disease mimics the valve problem, but which treatment fits the individual:

\begin{itemize}
    \item \textbf{Surgical replacement (SAVR)} versus \textbf{transcatheter replacement (TAVI)}
    \item \textbf{Medical management alone:} for people who decline any procedure or are too unwell to benefit
    \item \textbf{Balloon valvuloplasty:} a temporary widening of the valve, occasionally used as a bridge when immediate surgery is too risky
    \item \textbf{Valve repair:} sometimes possible in aortic regurgitation, preserving the native valve
    \item \textbf{Postoperative problems versus normal recovery:} wound infection, pneumonia, atrial fibrillation, or bleeding must be distinguished from the usual discomforts of recovery
\end{itemize}

\section*{When to Seek Medical Care}
Contact the cardiac team after surgery if you develop a fever, spreading redness or discharge from the wound, a wound that reopens, or pain that is worsening rather than settling. Report any new breathlessness, chest pain, fainting, or a fast, irregular heartbeat promptly, and attend all scheduled follow-up appointments and blood tests.

\textbf{Contact your local emergency services for:}
\begin{itemize}
    \item Severe chest pain or sudden breathlessness
    \item Fainting, collapse, or near-syncope
    \item Heavy bleeding from a wound
    \item Sudden weakness, numbness, facial droop, or difficulty speaking
\end{itemize}

\section*{Diagnosis and Investigations}
\begin{itemize}
    \item \textbf{Echocardiography:} the key test before surgery to assess the valve, the pressure gradient, and left ventricular function, and again afterwards to confirm the new valve is working well
    \item \textbf{Coronary angiography or cardiac CT:} to check the coronary arteries and plan the surgical or catheter approach, including assessment of the femoral arteries for TAVI
    \item \textbf{ECG and rhythm monitoring:} to detect and treat heart rhythm problems before and after the procedure
    \item \textbf{Blood tests:} kidney function, full blood count, and tests that guide medicines, including anticoagulation
    \item \textbf{Dental review:} treatment of tooth and gum disease before surgery reduces the risk of postoperative valve infection
    \item \textbf{Assessment of fitness:} a multidisciplinary evaluation of age, frailty, co-morbidities, and anatomy to choose the safest and most effective procedure
\end{itemize}

\section*{Management}
\begin{itemize}
    \item \textbf{Valve choice:} mechanical valves are very durable but require lifelong anticoagulation with warfarin; tissue (biological) valves usually require no long-term blood thinners but degenerate over roughly 10 to 20 years
    \item \textbf{Operative approach:} SAVR through a midline or mini chest incision, or TAVI through a catheter, usually via the groin and sometimes the subclavian artery or apex of the heart
    \item \textbf{Medicines after surgery:} blood thinners as needed, and medication for blood pressure, heart rhythm, or heart failure
    \item \textbf{Cardiac rehabilitation:} a supervised programme of graduated exercise and education that rebuilds fitness and confidence
    \item \textbf{Wound care:} keeping the wound clean and dry, and reporting any signs of infection promptly
    \item \textbf{Infection precautions:} excellent dental hygiene, prompt treatment of infections, and antibiotics before certain dental procedures where recommended
    \item \textbf{Lifelong follow-up:} regular cardiology review and echocardiography to monitor the valve for degeneration, leakage, or infection
\end{itemize}

\section*{Complications}
\begin{itemize}
    \item Bleeding, including around the heart (pericardial effusion or tamponade)
    \item Infection of the wound or of the new valve (endocarditis)
    \item Stroke or other thromboembolic complications
    \item Heart rhythm disturbances, most commonly atrial fibrillation
    \item The need for a permanent pacemaker, more frequent after TAVI
    \item Leakage around the new valve (paravalvular leak)
    \item With mechanical valves: bleeding risk from lifelong anticoagulation
    \item With tissue valves: gradual degeneration, eventually requiring further intervention
    \item Kidney injury, lung complications, and prolonged hospital stay, especially in frail patients
\end{itemize}

\section*{Prognosis and Outcomes}
Most people experience a dramatic improvement in breathlessness, chest pain, and quality of life after aortic valve replacement, and long-term survival after a successful procedure is good. Recovery typically takes several weeks to a few months, with activity rebuilt gradually through cardiac rehabilitation.

Mechanical valves offer excellent durability but carry the lifelong burden of anticoagulation and its bleeding risk; tissue valves avoid this but eventually wear out, which matters most in younger patients. TAVI has excellent outcomes in older and frailer patients, with a short hospital stay and rapid recovery. Individual prognosis depends on age, other health problems, the state of the heart before surgery, and how smoothly the operation and recovery run.

\section*{Prevention and Self-Care}
\begin{itemize}
    \item Stop smoking and control blood pressure, cholesterol, and diabetes to protect the remaining blood vessels
    \item Attend cardiac rehabilitation and build up physical activity gradually as advised
    \item Keep surgical wounds clean and dry, and watch for signs of infection
    \item Practise excellent dental hygiene, and follow advice about antibiotics before dental procedures
    \item Take all medicines exactly as prescribed, including blood thinners for mechanical valves, and attend monitoring appointments
    \item Carry information about your valve, and tell every health professional you see
    \item Report fevers, chills, or unusual symptoms promptly -- endocarditis must be caught early
    \item Maintain a heart-healthy diet, healthy weight, and regular physical activity within advised limits
\end{itemize}

\section*{Sources and Further Reading}
The following organisations provide reliable, up-to-date information on aortic valve replacement for patients, students, and health professionals.

\begin{itemize}
    \item NHS. \textit{Information on aortic valve replacement and recovery.} https://www.nhs.uk/conditions/
    \item NICE. \textit{Guidance on transcatheter aortic valve implantation and heart valve surgery.} https://www.nice.org.uk/
    \item Mayo Clinic. \textit{Overview of aortic valve surgery and valve replacement options.} https://www.mayoclinic.org/
    \item MedlinePlus. \textit{Health information on aortic valve surgery and aftercare.} https://medlineplus.gov/
    \item MSD Manual. \textit{Professional overview of heart valve surgery and prostheses.} https://www.msdmanuals.com/
    \item World Health Organization. \textit{Resources on cardiovascular disease and surgical care.} https://www.who.int/
\end{itemize}
